\section{Introduction}

Artificial Intelligence (AI) systems and their applications have been of significant consequence over recent decades. This has been especially evident in recent years abnormally significant awareness of AI in the general public. In particular, these include techniques from machine learning (ML) like the artificial neural networks (ANNs) that power contemporary generative AI (GEN-AI) systems like large language models (LLMs).

These AI systems have shown considerable promise, but the underlying technologies come with significant drawbacks. Concerns include the growing footprint of data centers, which has garnered public interest, the worsening use of finite resources such as fresh water and electricity, and the potential for substantial environmental and health impacts \parencite{Han2025,Siddik2021,Li2023}.
% Many ACM-formatted documents prefer fewer citations, at the end of the
% sentence. Moving things there just in case.

Conversely, Answer Set Programming (ASP) languages live at a different end of the research spectrum in terms of their efficiency. They are designed to solve NP-Complete and NP-Hard problems in a fast and efficient manner on limited hardware \parencite{Gebser2019}. Our current work involves exploring the use of ASP languages to accelerate the processes of learning and inference. The eventual goal of our research is to reduce the resources needed for these technologies.

In this paper we introduce our ANN representations in the Cling language. We also introduce the {\tt C-MCPN}: The \emph{Clingo McCulloch-Pitts Network Editor}. The C-MCPN is a tool that we developed to assist in the exploration of ASP-based networks. It allows a user to model the network architecture and conduct inference, with underlying representations and inference guided by declarative software written in Clingo. Specifically we restrict our attention to McCulloch-Pitts Networks (MPNs) as the target ANN type \parencite{McCulloch1943}.
